\section{From Alexandrov comparison to the tensor bound}
\label{sec:forward}

\begin{theorem}
\label{fwd:curvature}
Let \((M,d)\) be a complete \(n\)-dimensional Alexandrov space with
curvature bounded below by \(\kappa\), where \(n\geq2\), and put
\(m=\mathcal H^n_d\). Then \((M,d,m)\) is
\(\RCD((n-1)\kappa,n)\), its measure has full support, and
\begin{equation}
\label{fwd:target}
 R(X,Y,Y,X)(f)
 \geq \kappa\int f\bigl(|X|^2|Y|^2-\langle X,Y\rangle^2\bigr)\dd m
\end{equation}
for every nonnegative \(f\in\Test(M)\) and every
\(X,Y\in\TestV(M)\), with the global test classes and curvature
convention of Theorem~\ref{main:equivalence}.
\end{theorem}

The curvature-dimension and full-support assertions follow from
the Alexandrov-to-\(\CD\) theorem and infinitesimal Hilbertianity
\cite{Petrunin2011,ZhangZhu2010,KMS2001}. The unreduced statement
for arbitrary \(\kappa\), including complete Alexandrov spaces
with boundary, is recorded directly in
\cite[Theorem~2.9 and Corollary~2.10]{KapovitchKettererSturm};
it requires only local finiteness of \(\mathcal H^n\), with no
compactness or assumption of finite total mass. Set
\(K=(n-1)\kappa\). We shall use the resulting RCD calculus throughout
this section. In particular, all covariant derivatives below are weak
covariant derivatives. To prove the tensor inequality, we compare
two energies associated with short flows in finitely many gradient
directions: the action of the radial flows and the optimal transport
action between their endpoint measures. Alexandrov comparison bounds
twice the averaged radial action minus the averaged pair action from
below. For the upper bound, a dual transport calculation first
produces an error measuring radial acceleration. Second-order
endpoint corrections remove that error, after which averaging
identifies the upper bound as curvature plus a nonnegative square.
Sobolev dual potentials chosen separately on finitely many
localization pieces make this remaining square arbitrarily small.
Positivity and tensoriality then extend the resulting gradient inequality to test
coefficients; local approximation and exhaustion recover the global
test classes. The successive limits are specified in
Remark~\ref{fwd:limit-order}.

\subsection{Conventions and a short-flow expansion}

All unlabelled spatial integrals in this section are with respect to
\(m\). Write
\[
 |U\wedge V|^2:=|U|^2|V|^2-\langle U,V\rangle^2.
\]
\[
 \mathcal G:=\{\nabla u:
 u\in\Test(M),\ \operatorname{supp}u\text{ compact},\
 \Delta u\in L^\infty(M)\}.
\]
Every member of \(\mathcal G\) is bounded, has bounded divergence,
and has covariant derivative in \(L^2\). It belongs to the global
\(\TestV(M)\) by Remark~\ref{pre:bare-gradients}. Throughout we use
the stronger test cutoffs of Lemma~\ref{pre:approximation}\textup{(i)}.

For compactly supported Lipschitz \(f\), we use
\(R(X,Y,Z,W)(f)\) to denote the same integrated expression
as in the statement, even when \(f\notin\Test(M)\).
Lemma~\ref{pre:tensor-continuity} gives this extension, its tensor
identities, and continuity for weights with common compact support
that converge uniformly with weakly convergent \(L^2\) gradients.

We choose a finite probability law on compact test potentials with
bounded Laplacian, invariant under \(u\mapsto-u\). Let \(A=\nabla u_A\)
and \(B=\nabla u_B\) be independent choices from this law. The notation
\(\mathbb E\) always refers to this finite averaging. The example to
keep in mind is the uniform law giving the four fields
\(U,-U,V,-V\), with \(U,V\in\mathcal G\). Sign symmetry will cancel
the odd terms in the energy expansion, while the mixed pairs retain
the sectional-curvature expression for \(U,V\). We work with a
general finite law because the calculation uses only independence
and sign symmetry. In particular,
\begin{equation}
\label{fwd:centering}
 \mathbb E u_A=0,\qquad \mathbb E A=0,\qquad
 \mathbb E|A-B|^2=2\mathbb E|A|^2.
\end{equation}
For a fixed compactly supported Lipschitz function \(f\geq0\), and a
constant linear combination \(C\) of the chosen gradients, put
\[
 L_C:=-\operatorname{div}(fC)
      =-\langle\nabla f,C\rangle-f\operatorname{div}C.
\]
This is a bounded function with compact support. The case
\(\int f=0\) is trivial, so we may assume that \(\int f>0\).

For each choice of \(A\), let \(S_A\) be a bounded compact test field
with bounded divergence. These fields provide the second-order
endpoint corrections: the acceleration of the leading gradient flow
is \(\nabla_AA\), so they will approach \(-\nabla_AA\) only
after the limit in \(h\). For \(0<h<1\), let \(F^h_{r,A}\),
\(0\leq r\leq1\), be the regular Lagrangian flow of
\begin{equation}
\label{fwd:prescribed-velocity}
 b^h_{r,A}=hA+h^2rS_A.
\end{equation}
Denote its transported density and endpoint measure by
\[
 (F^h_{r,A})_\#(fm)=\rho^h_{r,A}m,
 \qquad \rho^h_A:=\rho^h_{1,A},\qquad \mu^h_A:=\rho^h_A m.
\]
All statements as \(h\downarrow0\) initially keep the finite family,
the weight, and the fields \(S_A\) fixed.

\begin{lemma}[Short-flow density expansion]
\label{fwd:flow}
The flows above exist and are invertible modulo null sets. Their
forward and inverse compression constants are \(1+O(h)\), their
displacements are \(O(h)\), and their transported densities have
support in a common compact set. Uniformly in \(r\in[0,1]\),
\begin{equation}
\label{fwd:flow-density-bound}
 \|\rho^h_{r,A}-f\|_{L^\infty}=O(h).
\end{equation}
Moreover,
\begin{equation}
\label{fwd:flow-first}
 \frac{\rho^h_{r,A}-f}{h}\rightharpoonup^* rL_A
\end{equation}
in \(L^\infty\) on the common bounded support, both at each fixed
\(r\) and on \([0,1]\times M\). For every fixed
\(q\in W^{1,2}(M)\),
\begin{equation}
\label{fwd:flow-second}
 \int q\rho^h_A
 =\int qf+h\int qL_A
 +\frac{h^2}{2}\int
   \langle\nabla q,fS_A+L_AA\rangle+o(h^2).
\end{equation}
If \(g\) is Lipschitz on a neighborhood of the common support, then
\begin{equation}
\label{fwd:increment}
 \frac{g(F^h_{1,A}(z))-g(z)}{h}
 \longrightarrow dg(A)(z)
 \quad\text{in }L^1(fm).
\end{equation}
\end{lemma}

\begin{proof}
The fields in \eqref{fwd:prescribed-velocity} are bounded, have
\(L^2\) covariant derivative and bounded divergence, and have support
in a fixed compact set. The Sobolev flow theorem applies
\cite[Theorems~4.3, 5.4, 6.7 and 8.3; Section~9.6]{AmbrosioTrevisan2014}. Applying it also to the time-reversed
field, and using uniqueness, gives the inverse flow. The divergence
bound is \(O(h)\), so both compression constants are \(1+O(h)\).
The speed is bounded by \(h\|A\|_\infty+h^2\|S_A\|_\infty\), which
gives the displacement and support assertions. Properness is used here
to replace bounded closed neighborhoods by compact sets.

To check \eqref{fwd:flow-density-bound}, the pushforward of \(m\)
under the flow has density between \(1-O(h)\) and \(1+O(h)\), by
the forward and inverse compression bounds. Replacing \(f(z)\) by
\(f(F^h_{r,A}(z))\) in the pushforward changes the density by at most
\(O(h)\operatorname{Lip}(f)\), again by compression and the displacement
bound. Multiplication of the pushed-forward reference density by
\(f\) changes it from \(f\) by at most \(O(h)\|f\|_\infty\).
This proves \eqref{fwd:flow-density-bound} without requiring a
spatially Lipschitz velocity field.

The continuity equation, tested with \(q\in W^{1,2}\), reads
\begin{equation}
\label{fwd:flow-ce}
 \int q(\rho^h_{r,A}-f)
 =\int_0^r\int
   \langle\nabla q,hA+h^2tS_A\rangle\rho^h_{t,A}\dd t.
\end{equation}
All pairings are finite because the densities are bounded and have
bounded support. Divide by \(h\) and use
\eqref{fwd:flow-density-bound}. The right-hand side converges,
uniformly in \(r\), to
\(r\int f\langle\nabla q,A\rangle=r\int qL_A\).
The uniform \(L^\infty\) bound on the left quotient, followed by
density of compact Lipschitz functions in \(L^1\) on bounded sets,
proves \eqref{fwd:flow-first}, also in space-time. Use
\eqref{fwd:flow-ce} once more, with \(r=1\), to obtain
\begin{align*}
 \int q(\rho^h_A-f)
 ={}&h\int f\langle\nabla q,A\rangle
 +h^2\int_0^1\int
   \left\langle\nabla q,
    \frac{\rho^h_{r,A}-f}{h}A+r\rho^h_{r,A}S_A\right\rangle\dd r.
\end{align*}
The functions paired with the normalized density difference are
fixed and integrable on its support. Equation
\eqref{fwd:flow-first} therefore gives \eqref{fwd:flow-second}.

For completeness, if \(v\in L^1\) on a common bounded neighborhood,
then
\[
 v\circ F^h_{r,A}\longrightarrow v
 \quad\text{in }L^1(\dd r\,fm).
\]
Indeed this holds for continuous functions by the displacement
bound, and follows for \(L^1\) functions by approximation and the
uniform compression bound. Differentiation of a Lipschitz function
along the flow gives
\[
 \frac{g(F^h_{1,A}(z))-g(z)}h
 =\int_0^1 dg(A)(F^h_{r,A}(z))\dd r
  +h\int_0^1r\,dg(S_A)(F^h_{r,A}(z))\dd r.
\]
The preceding composition convergence proves
\eqref{fwd:increment}. A cutoff may be inserted when \(g\) is only
locally Lipschitz.
\end{proof}

\subsection{The configuration comparison}

If \(\mu,\nu\) are finite nonnegative measures of equal mass, define
their transport energy by
\[
 J(\mu,\nu):=\frac12
 \inf_{\gamma\in\Pi(\mu,\nu)}\int d(x,y)^2\dd\gamma(x,y).
\]
This is the usual half squared Wasserstein distance after probability
normalization, multiplied by the common mass. Let
\[
 J^h_{AB}:=J(\mu^h_A,\mu^h_B),\qquad
 J^h_{0A}:=\frac12\int_0^1\int
       |hA+h^2rS_A|^2\rho^h_{r,A}\dd r.
\]
The second quantity is the energy of the prescribed radial flow,
not its optimal transport energy.

\begin{lemma}[Fourth-order configuration inequality]
\label{fwd:configuration}
For the above centered finite law and flows, put
\[
 a^h_A(z):=d(z,F^h_{1,A}(z)),\qquad
 e^h_{AB}(z):=d(F^h_{1,A}(z),F^h_{1,B}(z)).
\]
Then
\begin{equation}
\label{fwd:distance-first}
 \frac{a^h_A}{h}\longrightarrow |A|,\qquad
 \frac{e^h_{AB}}h\longrightarrow |A-B|
\end{equation}
in measure and in every finite \(L^p(fm)\), for each choice of
\(A,B\). Furthermore,
\begin{equation}
\label{fwd:energy-comparison}
 \liminf_{h\downarrow0}
 \frac{2\mathbb E J^h_{0A}-\mathbb E J^h_{AB}}{h^4}
 \geq\frac\kappa6\int f\,\mathbb E|A\wedge B|^2.
\end{equation}
\end{lemma}

\begin{proof}
At \(m\)-almost every base point \(z\), the Alexandrov tangent cone
is Euclidean. Choose a geodesic from \(z\) to each endpoint
\(F^h_{1,A}(z)\). For the finite weights \(p_A\), the radial vectors
\(v_A\in T_zM\) satisfy
\[
 \sum_{A,B}p_Ap_B|v_A||v_B|\cos\angle(v_A,v_B)
 =\left|\sum_Ap_Av_A\right|^2\geq0.
\]
Actual angles are at least the corresponding \(\kappa\)-comparison
angles. Consequently,
\begin{equation}
\label{fwd:cosine-sum}
 \mathbb E\!\left[
 a^h_Aa^h_B\cos\widetilde\angle_\kappa
          (a^h_A,a^h_B,e^h_{AB})\right]\geq0.
\end{equation}
Zero radial lengths contribute zero. All sides are \(O(h)\), so the
model triangles are defined for sufficiently small \(h\), including
when \(\kappa>0\).

For radial side lengths \(a,b\) and opposite side length \(e\), put
\(z_0=(a^2+b^2-e^2)/2\). The model cosine law gives, uniformly for
such small triangles,
\begin{equation}
\label{fwd:model-expansion}
 ab\cos\widetilde\angle_\kappa
 =z_0-\frac\kappa6(a^2b^2-z_0^2)+O(h^6).
\end{equation}
For \(\kappa=0\) this is exact without the correction. For
\(\kappa\ne0\), substitute the Taylor series of the model sine and
cosine into their cosine law. After multiplication by \(ab\), the
denominators are the nonvanishing analytic factors
\(\mathfrak s_\kappa(a)/a\) and
\(\mathfrak s_\kappa(b)/b\), where
\(\mathfrak s_\kappa''+\kappa\mathfrak s_\kappa=0\),
\(\mathfrak s_\kappa(0)=0\), and
\(\mathfrak s_\kappa'(0)=1\). This also proves uniformity at
degenerate triangles and zero sides.

We first identify the leading distances. By integration of the speed
and the composition convergence used in Lemma~\ref{fwd:flow},
\[
 \left(\frac{a^h_A}{h}-|A|\right)_+
 \longrightarrow0\quad\text{in }L^1(fm).
\]
There is a countable collection of scalar \(1\)-Lipschitz functions
whose differentials norm every tangent vector almost everywhere;
one may use the norming family in
\cite[Lemma 9.2]{AmbrosioTrevisan2014}, with local cutoffs when needed.
For each member \(g\), \eqref{fwd:increment} and its difference for
two flows give the lower bounds
\[
 a^h_A/h\geq
 |g(F^h_{1,A})-g|/h,\qquad
 e^h_{AB}/h\geq
 |g(F^h_{1,A})-g(F^h_{1,B})|/h.
\]
First taking a finite part of the norming family, and then increasing
it, shows that the negative parts of
\(a^h_A/h-|A|\) and \(e^h_{AB}/h-|A-B|\) tend to zero in measure.
The radial convergence follows. All scaled distances are uniformly
bounded, so this convergence also holds in every finite \(L^p(fm)\).

Equations \eqref{fwd:cosine-sum}--\eqref{fwd:model-expansion}, initially
retaining only their quadratic terms, imply
\[
 \mathbb E\int f(e^h_{AB}/h)^2
 \leq 2\mathbb E\int f(a^h_A/h)^2+O(h^2)
 \longrightarrow 2\mathbb E\int f|A|^2.
\]
By \eqref{fwd:centering}, the last quantity is
\(\mathbb E\int f|A-B|^2\). The pair-distance lower bounds already
proved imply convergence of the negative errors in \(L^1\).
The displayed upper bound therefore forces the positive errors to
converge in \(L^1\) as well. Since the family is finite, this proves
the pair convergence in \eqref{fwd:distance-first} for every pair of
positive probability.

We may now retain the quartic term in
\eqref{fwd:model-expansion}. The convergence just proved gives
\[
 h^{-4}\left((a^h_Aa^h_B)^2
 -\left(\frac{(a^h_A)^2+(a^h_B)^2-(e^h_{AB})^2}{2}\right)^2\right)
 \longrightarrow |A\wedge B|^2
\]
in \(L^1(fm)\). Thus
\begin{equation}
\label{fwd:distance-fourth}
 \liminf_{h\downarrow0}\frac1{h^4}
 \left(\mathbb E\int f(a^h_A)^2
       -\frac12\mathbb E\int f(e^h_{AB})^2\right)
 \geq\frac\kappa6\int f\,\mathbb E|A\wedge B|^2.
\end{equation}
The radial energy dominates half the squared radial displacement,
and \((F^h_{1,A},F^h_{1,B})_\#(fm)\) is a coupling of the endpoint
measures. Hence
\[
 2J^h_{0A}\geq\int f(a^h_A)^2,
 \qquad J^h_{AB}\leq\frac12\int f(e^h_{AB})^2.
\]
Their combination with \eqref{fwd:distance-fourth} proves
\eqref{fwd:energy-comparison}.
\end{proof}

\subsection{First variation of the optimal interpolation densities}

The transport expansion will pair the first density variation with
fixed \(L^2\) Hessian coefficients. We therefore need an \(O(h)\)
bound in \(L^2\), as well as identification of the weak limit.
For the next lemma, assume that \(f\) is semiconcave with some finite
constant along geodesics in a bounded region containing all the
interpolations under consideration. No bound for that constant is
required uniformly over later approximations of \(f\).
Let \(\sigma^h_s m\), \(0\leq s\leq1\), be the optimal interpolation
from \(\mu^h_A\) to \(\mu^h_B\), and let \(v^h_s\) be its minimal
continuity-equation velocity. Thus
\[
 2J^h_{AB}=\int_0^1\int\sigma^h_s|v^h_s|^2\dd s.
\]
Existence, uniqueness, and bounded densities are used in their usual
finite-dimensional RCD form
\cite{GigliRajalaSturm2016,Rajala2012}. We derive the needed
pointwise distortion inequality from integrated \(\CD\) by
restriction below.

\begin{lemma}[Optimal density first variation]
\label{fwd:density}
Put \(D=B-A\) and \(C_s=(1-s)A+sB\). The densities \(\sigma^h_s\)
have common compact support and a uniform \(L^\infty\) bound. They
satisfy
\begin{equation}
\label{fwd:density-bound}
 \sup_{0\leq s\leq1}\|\sigma^h_s-f\|_{L^2}^2=O(h^2),
 \qquad
 \frac{\sigma^h_s-f}{h}\rightharpoonup L_{C_s}
 \quad\text{in }L^2(\dd s\dd m).
\end{equation}
Their velocities satisfy
\begin{equation}
\label{fwd:velocity-first}
 \int_0^1\int\sigma^h_s|v^h_s/h-D|^2\dd s\longrightarrow0.
\end{equation}
\end{lemma}

\begin{proof}
\emph{Restriction of the plan.}
The endpoints have uniformly bounded densities and support in a
common ball. All connecting geodesics lie in a larger fixed ball,
which is compact. The standard interpolation bound gives the uniform
density estimate. We record a sharper consequence needed below.
Weaken the Ricci parameter to a fixed \(K'<0\). The zero-mass case
is immediate; otherwise let \(\pi^h\) be the optimal dynamical
plan and normalize it to a probability \(\pi\), with marginals
\(r_t m\). By
\cite[Theorem~1.1]{GigliRajalaSturm2016}, \(\pi\) has a Borel
conull set on which every evaluation \(e_t\), \(t<1\), is
injective. Both endpoints are absolutely continuous. Applying the
same theorem to the reversed plan, pulling its conull set back
under time reversal, and intersecting with the first set gives a
Borel conull \(\Gamma\) on which \(e_t\) is
injective for every \(t\in[0,1]\).

Fix \(s\in(0,1)\), and let \(H\geq0\) be bounded Borel on
\(\Gamma\), with \(Z=\int H\dd\pi>0\). The probability
\(\pi_H=(H/Z)\pi\) is optimal for its endpoint marginals:
after scaling \(H\leq1\), a cheaper transport for \(H\pi\)
could be spliced with \(\pi-H\pi\), contradicting optimality.
Its endpoint densities are bounded and absolutely continuous.
The integrated \(\CD(K',n)\) inequality supplies an optimal
dynamical plan for these marginals, and uniqueness forces that plan
to be \(\pi_H\). The measurable inverses of \(e_t|_\Gamma\)
give, by Lusin--Souslin,
\[
 h_t=H\circ(e_t|_\Gamma)^{-1},\qquad
 r_t^H=(h_t/Z)r_t,\qquad t\in\{0,s,1\},
\]
where \(h_t\) is extended by zero off \(e_t\Gamma\).
Write \(\beta_t(\ell)=\tau_{K',n}^{(t)}(\ell)\) for the
unreduced distortion coefficient and set
\[
 F_s(\gamma)=r_s(\gamma_s)^{-1/n}
 -\beta_{1-s}(\ell)r_0(\gamma_0)^{-1/n}
 -\beta_s(\ell)r_1(\gamma_1)^{-1/n}.
\]
The integrated inequality for \(\pi_H\), after cancellation of
\(Z^{1/n-1}\), is exactly
\[
 \int H^{1-1/n}F_s\dd\pi\geq0.
\]
There are no undefined inverse powers in this calculation: zero or
infinite density values lie over marginal-null sets at the three
fixed times, and on the common finite-measure compact support \(Q\),
\[
 \int r_t(\gamma_t)^{-1/n}\dd\pi
 =\int_Q r_t^{1-1/n}\dd m<\infty.
\]
Taking \(H\) to be indicators of Borel sets proves
\(F_s\geq0\) \(\pi\)-almost everywhere. Restoring the common
mass preserves this homogeneous inequality. Its exceptional set
may depend on \(s\); every fixed-time estimate below has the same
constants, which suffices for the supremum over \(s\).
Jointly Borel density versions and Fubini give the corresponding
\(\dd s\otimes\pi\)-almost-everywhere statements for time integrals.

\emph{The density norm bound.}
On the bounded range of possible geodesic lengths \(\ell\), the
pointwise CD distortion coefficients, divided by their zero-curvature
coefficients, have a positive lower bound and are at least
\(1-C\ell^2\), uniformly in \(s\). The negative power mean is at
most the arithmetic mean. Since the endpoint densities are bounded,
the pointwise CD inequality therefore implies, along the optimal
dynamical plan \(\pi^h\),
\begin{equation}
\label{fwd:pointwise-density}
 \sigma^h_s(\gamma_s)
 \leq (1-s)\rho^h_A(\gamma_0)
       +s\rho^h_B(\gamma_1)+C\ell(\gamma)^2.
\end{equation}
To see that the error remains valid when \(\ell\) is not small,
let \(a(\ell)\) be a common positive lower bound for the two
distortion ratios. The CD inequality first gives the negative power
mean multiplied by \(a(\ell)^{-n}\). On the fixed length interval,
\(a(\ell)^{-n}-1\leq C\ell^2\). This proves
\eqref{fwd:pointwise-density}, including such lengths.

Integration against \(\pi^h\) gives
\[
 \int(\sigma^h_s)^2
 \leq(1-s)\int(\rho^h_A)^2+s\int(\rho^h_B)^2
      +C\int\ell^2\dd\pi^h.
\]
Semiconcavity of \(f\) gives
\[
 -2\int f\sigma^h_s
 \leq-2(1-s)\int f\rho^h_A-2s\int f\rho^h_B
      +C\int\ell^2\dd\pi^h.
\]
Here \(\int\ell^2\dd\pi^h=2J^h_{AB}=O(h^2)\), by the
configuration coupling. Adding the last two inequalities and
\(\int f^2\) proves the norm bound in
\eqref{fwd:density-bound}, using
\eqref{fwd:flow-density-bound} for the endpoints.

\emph{First variations of velocity and density.}
The same coupling and \eqref{fwd:distance-first} imply
\[
 \limsup_{h\downarrow0}\int_0^1\int
     \sigma^h_s|v^h_s/h|^2\dd s\leq\int f|D|^2.
\]
The continuity equation tested with \(u_B-u_A\), and then
\eqref{fwd:flow-second}, gives
\[
 \int_0^1\int\sigma^h_s\langle v^h_s/h,D\rangle\dd s
 =\frac1h\int(u_B-u_A)(\rho^h_B-\rho^h_A)
 \longrightarrow\int f|D|^2.
\]
The norm estimate for the densities also gives
\(\int_0^1\int\sigma^h_s|D|^2\dd s\to\int f|D|^2\).
Expanding the square proves \eqref{fwd:velocity-first}.

Finally, for any fixed \(q\in W^{1,2}(M)\),
\[
 \int q\frac{\sigma^h_s-f}{h}
 =\int q\frac{\rho^h_A-f}{h}
   +\int_0^s\int\sigma^h_t\langle\nabla q,v^h_t/h\rangle\dd t.
\]
The first term converges to \(\int qL_A\). In the second, use
\eqref{fwd:velocity-first} and Cauchy--Schwarz to replace
\(v^h_t/h\) by \(D\); then use the density norm bound to replace
\(\sigma^h_t\) by \(f\). The error tends to zero uniformly in
\(s\). The limit is \(s\int f\langle\nabla q,D\rangle
=s\int qL_D\). Boundedness of the normalized density differences
in \(L^2(\dd s\dd m)\), and density of finite sums of separated
space-time test functions, prove the weak convergence in
\eqref{fwd:density-bound}.
\end{proof}

\subsection{The fourth-order transport estimate}

The configuration comparison gives a lower bound for
\(2\mathbb E J^h_{0A}-\mathbb E J^h_{AB}\). We now obtain an
upper bound for the same difference in terms of curvature and a
nonnegative square, which the following subsection removes by
localization.

Continue to assume that \(f\) is semiconcave as above, keeping the
successive limits of Remark~\ref{fwd:limit-order}. For each
ordered pair \(A,B\), choose an arbitrary fixed
\(q_{AB}\in W^{1,2}(M)\). The choices need not have any symmetry.
These are dual energy tests. After extending the error estimate to
Lipschitz weights, we will choose their gradients on finitely many
localization pieces to approximate \(-\nabla_A(B-A)\) and remove
the remaining square. Each potential is held fixed during the
short-flow limit.
For the duration of a calculation with this pair write
\[
 D=B-A,\qquad C_s=(1-s)A+sB,\qquad T_D=\nabla_DD,
\]
\[
 q_s=q_{AB}-s\frac{|D|^2}{2},\qquad
 Q_s=\nabla q_s=\nabla q_{AB}-sT_D.
\]
The identity for \(Q_s\) follows from Hessian symmetry, since
\(D\) is a gradient. Notice that \(T_D\in L^2\) and that
\(|D|^2/2\in W^{1,2}\). No Hessian of \(q_{AB}\) is required.

\begin{lemma}[Dual expansion and radial defect]
\label{fwd:dual}
With the preceding notation,
\begin{align}
 J^h_{AB}\geq{}&
 h\int(u_B-u_A)(\rho^h_B-\rho^h_A)
 -\frac{h^2}{2}\int f|D|^2
 -h^3\int_0^1\int f\langle C_s,T_D\rangle\dd s
 \nonumber\\
 &\quad+h^4K_{AB}(S)+o(h^4),
 \label{fwd:dual-expansion}\\
 K_{AB}(S):={}&
 \frac12\int\left[
  \langle Q_1,fS_B+L_BB\rangle
  -\langle Q_0,fS_A+L_AA\rangle\right]
 \nonumber\\
 &\quad-\int_0^1\int
     \left[L_{C_s}\langle D,Q_s\rangle+\frac f2|Q_s|^2\right]\dd s.
 \label{fwd:dual-coefficient}
\end{align}
For the prescribed radial action \(J^h_{0A}\), take pair indices
\(0,A\), zero initial potential and fields, \(\rho^h_0=f\), and
\(q_{0A}=0\). Then \eqref{fwd:dual-expansion} holds with equality
and fourth-order coefficient
\begin{equation}
\label{fwd:radial-defect}
 K_{0A}(S)+\frac16\int f|S_A+\nabla_AA|^2.
\end{equation}
\end{lemma}

\begin{proof}
\emph{The energy identity.}
For a bounded-compression solution \(\sigma_s\) of the continuity
equation with velocity \(v_s\), use
\[
 p_s=h(u_B-u_A)+h^2q_s.
\]
Its strong \(W^{1,2}\) dependence on time and the continuity
equation give the exact energy identity
\begin{align}
 \frac12\int_0^1\int\sigma_s|v_s|^2\dd s
 ={}&\left[\int p_s\sigma_s\right]_{s=0}^{s=1}
 -\int_0^1\int\sigma_s
      \left(\partial_sp_s+\frac12|\nabla p_s|^2\right)\dd s
 \nonumber\\
 &\quad+\frac12\int_0^1\int\sigma_s|v_s-\nabla p_s|^2\dd s.
 \label{fwd:young-identity}
\end{align}
One may verify the time differentiation directly by writing \(p_s\)
as a sum of fixed Sobolev functions with affine time coefficients.
The endpoint terms are finite by the bounded densities and bounded
support; the remaining terms are finite by the kinetic-energy and
\(W^{1,2}\) bounds. Thus no boundedness assumption on \(q_{AB}\)
is needed.

\emph{The dual expansion.}
The term \(-s|D|^2/2\) in \(q_s\) cancels the quadratic part of
the Hamilton--Jacobi residual. Thus the residual begins at order
\(h^3\), and the first density variation supplies its contribution
at order \(h^4\). Explicitly,
\begin{equation}
\label{fwd:residual}
 \partial_sp_s+\frac12|\nabla p_s|^2
 =h^3\langle D,Q_s\rangle+\frac{h^4}{2}|Q_s|^2.
\end{equation}
For the optimal interpolation, Lemma~\ref{fwd:density} gives
\begin{align*}
 \int_0^1\int\sigma^h_s\langle D,Q_s\rangle\dd s
 ={}&\int_0^1\int f\langle D,Q_s\rangle\dd s
 +h\int_0^1\int L_{C_s}\langle D,Q_s\rangle\dd s+o(h),\\
 \int_0^1\int\sigma^h_s|Q_s|^2\dd s
 ={}&\int_0^1\int f|Q_s|^2\dd s+o(1).
\end{align*}
The first limit uses weak \(L^2\) convergence against the fixed
\(L^2\) function \(\langle D,Q_s\rangle\). The second uses convergence
in measure and the uniform density bound, by truncation of the fixed
integrable function \(|Q_s|^2\). Outside the common density support
both statements are vacuous.

Apply \eqref{fwd:flow-second} to the two fixed endpoint functions
\(q_0,q_1\). Leave the endpoint contribution involving
\(u_B-u_A\) unexpanded. The terms of order four are then precisely
\eqref{fwd:dual-coefficient}. For the cubic term, integration by
parts gives
\[
 \int(q_1L_B-q_0L_A)
 =\int f(\langle Q_1,B\rangle-\langle Q_0,A\rangle).
\]
Using \(Q_1=Q_0-T_D\) and
\(\int_0^1Q_s\dd s=Q_0-T_D/2\), its difference from the cubic
residual is
\[
 -\int f\left\langle B-\frac D2,T_D\right\rangle
 =-\int_0^1\int f\langle C_s,T_D\rangle\dd s.
\]
Also \(q_1-q_0=-|D|^2/2\). Dropping the nonnegative last term in
\eqref{fwd:young-identity} proves \eqref{fwd:dual-expansion}.

\emph{The radial defect.}
For the radial flow use its actual velocity from
\eqref{fwd:prescribed-velocity}, and use Lemma~\ref{fwd:flow} in
place of Lemma~\ref{fwd:density}. Here \(D=A\), \(C_s=sA\), and
\(Q_s=-s\nabla_AA\). The last term of
\eqref{fwd:young-identity}, divided by \(h^4\), is
\[
 \frac12\int_0^1s^2\int
     \rho^h_{s,A}|S_A+\nabla_AA|^2\dd s
 \longrightarrow\frac16\int f|S_A+\nabla_AA|^2.
\]
The integrand in the spatial limit is fixed and integrable, and
\eqref{fwd:flow-density-bound} justifies the limit. This proves the
radial assertion.
\end{proof}

We next compute the fourth-order difference. Define
\[
 F_C:=L_CC-f\nabla_CC
\]
and the separately quadratic expression
\begin{equation}
\label{fwd:P-definition}
 P(C,C;D,D):=\int\left[
       -\langle F_C,\nabla_DD\rangle+f|\nabla_CD|^2\right].
\end{equation}
The notation \(P(E,F;G,H)\) means its polarization, symmetric in
the first pair and separately symmetric in the second pair. Let
\(K_{AB}\) denote the expression in \eqref{fwd:dual-coefficient}
with \(fS_E+L_EE\) replaced by \(F_E\) at both endpoints.

\begin{lemma}[Square completion and curvature identity]
\label{fwd:algebra}
For every ordered pair,
\begin{equation}
\label{fwd:square-completion}
 K_{AB}=\frac12\int_0^1
 \left[P(C_s,C_s;D,D)-\int f|Q_s+\nabla_{C_s}D|^2\right]\dd s.
\end{equation}
For the radial coefficients use \(q_{0A}=0\). Then
\begin{align}
 2\mathbb E K_{0A}-\mathbb E K_{AB}
 ={}&\frac13\mathbb E
       [P(A,B;A,B)-P(A,A;B,B)]
 \nonumber\\
 &\quad+\frac12\mathbb E\int
 f|\nabla q_{AB}+\nabla_A(B-A)|^2
 \label{fwd:average-P}\\
 ={}&\frac16\mathbb ER(A,B,B,A)(f)
 +\frac12\mathbb E\int
 f|\nabla q_{AB}+\nabla_A(B-A)|^2.
 \label{fwd:average-curvature}
\end{align}
\end{lemma}

\begin{proof}
\emph{Integration by parts.}
We first justify the identity
\begin{equation}
\label{fwd:gradient-ibp}
 \int\langle L_DC-L_CD,\nabla q\rangle
 =\int f\langle\nabla_DC-\nabla_CD,\nabla q\rangle,
 \qquad q\in W^{1,2}(M).
\end{equation}
For test \(q\), integration by parts in its two terms gives
\begin{align*}
 \int L_D\langle C,\nabla q\rangle
 &=\int f\bigl(\langle\nabla_DC,\nabla q\rangle
                    +\Hess q(D,C)\bigr),\\
 \int L_C\langle D,\nabla q\rangle
 &=\int f\bigl(\langle\nabla_CD,\nabla q\rangle
                    +\Hess q(C,D)\bigr).
\end{align*}
The Hessian terms cancel. The vector fields paired with
\(\nabla q\) on both sides of \eqref{fwd:gradient-ibp} are in
\(L^2\) and have bounded support. Density of test functions in
\(W^{1,2}\) therefore proves the identity for the stated class,
without assigning a Hessian to \(q\).

\emph{Square completion.}
The maps \(s\mapsto F_{C_s}\) and \(s\mapsto Q_s\) are
polynomials with values in \(L^2\). They have derivatives
\[
 F_{C_s}'=L_DC_s+L_{C_s}D-f\nabla_DC_s-f\nabla_{C_s}D,
 \qquad Q_s'=-\nabla_DD.
\]
Express the endpoint difference in \(K_{AB}\) as the integral of
the derivative of \(\langle F_{C_s},Q_s\rangle/2\), and subtract
the other terms in \eqref{fwd:dual-coefficient}. Applying
\eqref{fwd:gradient-ibp} gives
\[
 K_{AB}=\frac12\int_0^1\int
 \left[-\langle F_{C_s},\nabla_DD\rangle
       -2f\langle\nabla_{C_s}D,Q_s\rangle-f|Q_s|^2\right]\dd s.
\]
Completion of the square proves \eqref{fwd:square-completion}.

\emph{Averaging.}
For the averaging, set
\[
 T=\mathbb E P(A,A;A,A),\quad
 S=\mathbb E P(A,A;B,B),\quad
 H=\mathbb E P(A,B;A,B).
\]
Independence and sign symmetry kill every term with an odd number
of occurrences of either random field. Exchangeability identifies
the two pure and the two unmixed terms. Expansion in the four slots
therefore gives
\[
 \mathbb E P(C_s,C_s;D,D)
 =((1-s)^2+s^2)(T+S)-4s(1-s)H.
\]
Both \(\int_0^1((1-s)^2+s^2)\dd s\) and
\(\int_0^14s(1-s)\dd s\) equal \(2/3\). The contribution of the
\(P\) term to \(\mathbb E K_{AB}\) is consequently
\((T+S-H)/3\). In the radial case \(C_s=sA\), \(D=A\),
\(Q_s=-s\nabla_AA\), so the square vanishes and
\(K_{0A}=P(A,A;A,A)/6\). Finally,
\[
 Q_s+\nabla_{C_s}D=\nabla q_{AB}+\nabla_AD.
\]
These facts prove \eqref{fwd:average-P}.

\emph{The curvature contraction.}
To identify the curvature convention explicitly, write
\[
 a=\nabla_AA,\qquad b=\nabla_BB,\qquad
 U=\nabla_BA,\qquad V=\nabla_AB.
\]
Separate symmetric polarization of \eqref{fwd:P-definition} gives
\begin{align}
 P(A,B;A,B)=\int\biggl[
 &-\frac14\langle L_AB+L_BA,U+V\rangle
 \nonumber\\
 &+f\left(\frac14|U|^2+\frac14|V|^2
             +\langle U,V\rangle+\frac12\langle a,b\rangle\right)
 \biggr],
 \label{fwd:polarization}
\end{align}
whereas
\[
 P(A,A;B,B)=\int
 [-L_A\langle A,b\rangle+f\langle a,b\rangle+f|V|^2].
\]
Exchangeability gives \(\mathbb E|U|^2=\mathbb E|V|^2\) and
\[
 \frac12\mathbb E\int\langle L_AB+L_BA,U+V\rangle
 =\mathbb E\int L_B\langle A,U+V\rangle.
\]
Hessian symmetry for the gradient \(A\) gives
\(\langle A,U\rangle=\langle B,a\rangle\); swapping the random
choices then yields
\(\mathbb E\int L_B\langle A,U\rangle
 =\mathbb E\int L_A\langle A,b\rangle\).
Thus
\begin{align}
 2\mathbb E[P(A,B;A,B)-P(A,A;B,B)]
 =\mathbb E\int\bigl[
 &L_A\langle A,b\rangle-L_B\langle A,V\rangle
 \nonumber\\
 &+f(-\langle a,b\rangle-|V|^2+2\langle U,V\rangle)
 \bigr].
 \label{fwd:polarized-curvature}
\end{align}
The first two terms of the specified integrated curvature formula,
with \((X,Y,Z,W)=(A,B,B,A)\), combine to
\[
 L_A\langle A,b\rangle-f\langle a,b\rangle.
\]
Its next two terms combine to
\[
 -L_B\langle A,V\rangle+f\langle U,V\rangle.
\]
Finally \([A,B]=V-U\), and Hessian symmetry for \(B\) gives
\[
 -f\langle\nabla_{[A,B]}B,A\rangle
 =-f\langle V-U,V\rangle.
\]
The sum is exactly the integrand in
\eqref{fwd:polarized-curvature}. This proves
\eqref{fwd:average-curvature}, including the sign of the bracket term.
\end{proof}

\begin{proposition}[A localized error estimate]
\label{fwd:defect}
For the finite centered law above, every semiconcave compactly
supported Lipschitz \(f\geq0\), and every family
\(q_{AB}\in W^{1,2}(M)\),
\begin{equation}
\label{fwd:defective-curvature}
 \frac\kappa6\int f\,\mathbb E|A\wedge B|^2
 \leq\frac16\mathbb ER(A,B,B,A)(f)
 +\frac12\mathbb E\int
 f|\nabla q_{AB}+\nabla_A(B-A)|^2.
\end{equation}
\end{proposition}

\begin{proof}
First keep all endpoint corrections \(S_A\) fixed. In twice the
expected radial expansion from Lemma~\ref{fwd:dual}, minus the
expected pair lower bound, every term displayed before order four
cancels. Indeed independence and \(\mathbb E u_A=0\) give the exact
identity
\[
 \mathbb E\int(u_B-u_A)(\rho^h_B-\rho^h_A)
 =2\mathbb E\int u_A(\rho^h_A-f).
\]
This identity uses no further expansion of the endpoint densities.
The quadratic terms cancel by \eqref{fwd:centering}. The radial
cubic term is a constant multiple of
\(\int f\langle A,\nabla_AA\rangle\), which is odd under sign
reversal. The pair cubic term is a constant multiple of
\[
 \int f\left\langle\frac{A+B}{2},\nabla_{B-A}(B-A)\right\rangle,
\]
which is odd under simultaneous sign reversal of both choices.
Their expectations therefore vanish. Together with
Lemma~\ref{fwd:configuration}, this proves
\begin{equation}
\label{fwd:fixed-correction-bound}
 \frac\kappa6\int f\,\mathbb E|A\wedge B|^2
 \leq 2\mathbb E K_{0A}(S)-\mathbb E K_{AB}(S)
       +\frac13\mathbb E\int f|S_A+\nabla_AA|^2.
\end{equation}

Now approximate \(-\nabla_AA\) in \(L^2\) by admissible fields
\(S_A^{(j)}\). Such an approximation can be chosen explicitly:
\(|A|^2/2\) is bounded, compactly supported, and Sobolev; heat
regularize it, multiply by a fixed test cutoff equal to one on a
neighborhood of its support, and take the negative gradient.
Strong continuity of the heat flow in \(W^{1,2}\) gives
\[
 S_A^{(j)}\longrightarrow-\nabla_AA\quad\text{in }L^2.
\]
Each approximant is a bounded compact test gradient with bounded
divergence. Its bounds may depend on \(j\). For every fixed \(j\),
the limit \(h\downarrow0\) leading to
\eqref{fwd:fixed-correction-bound} has already been taken, so no
uniform flow bounds in \(j\) are required.

The dependence of \(K_{AB}(S)\) on an endpoint correction is a
pairing with one of the fixed \(L^2\) fields \(fQ_0,fQ_1\).
It is therefore continuous under this approximation. The radial
defects tend to zero. Taking \(j\to\infty\) in
\eqref{fwd:fixed-correction-bound} gives
\[
 \frac\kappa6\int f\,\mathbb E|A\wedge B|^2
 \leq2\mathbb E K_{0A}-\mathbb E K_{AB}.
\]
Lemma~\ref{fwd:algebra} is exactly
\eqref{fwd:defective-curvature}.
\end{proof}

\subsection{Removing the error by localization}

\begin{lemma}[Semiconcave approximation of the weight]
\label{fwd:weight-approximation}
If \(f\geq0\) is compactly supported and Lipschitz, then there are
nonnegative compactly supported semiconcave functions \(f_j\), with
support in \(\operatorname{supp}f\), such that \(f_j\to f\)
uniformly and weakly in \(W^{1,2}\), and
\(\sup_j\operatorname{Lip}(f_j)<\infty\). Consequently
\eqref{fwd:defective-curvature} holds for every nonnegative compactly
supported Lipschitz weight.
\end{lemma}

\begin{proof}
Use the inf-convolutions
\[
 f_\varepsilon(x)
 :=\inf_y\left(f(y)+\frac{d(x,y)^2}{2\varepsilon}\right),
 \qquad\varepsilon\downarrow0.
\]
Minimizers exist by properness. If \(L=\operatorname{Lip}(f)\), then
\[
 0\leq f_\varepsilon\leq f,
 \qquad \|f_\varepsilon-f\|_\infty\leq\frac{\varepsilon L^2}{2}.
\]
In particular the supports lie in the fixed compact support of \(f\).
Every minimizer \(y\) satisfies \(d(x,y)\leq\varepsilon L\):
compare it with a point a short distance along a minimizing segment
from \(y\) to \(x\), divide the resulting optimality inequality by
that distance, and let the distance tend to zero. The same bound
in the difference-quotient estimate for the inf-convolution gives
local slope at most \(L\); since \(M\) is a length space,
\(\operatorname{Lip}(f_\varepsilon)\leq L\).

For sufficiently small fixed \(\varepsilon\), all these minimizing
distances are below a fixed model-comparison scale. Triangle
comparison makes the functions \(x\mapsto d(x,y)^2\) uniformly
semiconcave when these distances remain below that scale, including
at \(x=y\); compare \cite[Section~1.2]{Petrunin2007}.
For an explicit normalization, choose \(a>0\) with
\(-a^2\le\kappa\), weaken the curvature bound to \(-a^2\), and put \(h(r)=(\cosh(ar)-1)/a^2\). Model comparison
gives \((h\circ d_y)''\le1+a^2h\circ d_y\) along unit-speed
geodesics. Write \(r^2=\Psi(h(r))\). The function \(\Psi\) is
increasing and concave, is smooth at zero, and satisfies
\(0\le\Psi'\le2\). Thus on \(d_y\le R\) the squared distance has
semiconcavity constant \(2\cosh(aR)\), uniformly in the center. Locally near each \(x\), the
infimum defining \(f_\varepsilon\) may be restricted to a common
small neighborhood of \(x\), by the minimizer bound. It is thus
an infimum of functions with a common semiconcavity constant there.
Infima preserve that semiconcavity inequality. The constant depends
on \(\varepsilon,\kappa,L\), but not on the base point. Local
semiconcavity with this constant yields the same inequality along
each geodesic by subdivision. This proves the required
semiconcavity of each \(f_\varepsilon\).

The common support and Lipschitz bound give boundedness in
\(W^{1,2}\). Uniform convergence identifies every weak Sobolev
subsequential limit with \(f\); hence any sequence
\(\varepsilon_j\downarrow0\) has the asserted weak convergence.
In \eqref{fwd:defective-curvature}, the two terms without curvature
are integrals against fixed locally integrable functions. The
curvature term converges by Lemma~\ref{pre:tensor-continuity}. Passing
to the limit proves the last assertion.
\end{proof}

\begin{lemma}[Approximation by finitely many gradient choices]
\label{fwd:gradient-selection}
Let \(f\geq0\) be compactly supported and Lipschitz, and let
\(V_1,\ldots,V_N\in L^2(TM)\). For every \(\eta>0\), there are
finitely many nonnegative Lipschitz functions
\(\alpha_1,\ldots,\alpha_L\), with
\(\sum_{\ell=1}^L\alpha_\ell=1\), and functions
\(q_{r,\ell}\in W^{1,2}(M)\), such that
\begin{equation}
\label{fwd:selection-error}
 \sum_{r=1}^N\sum_{\ell=1}^L
 \int f\alpha_\ell|\nabla q_{r,\ell}-V_r|^2<\eta.
\end{equation}
The products \(f\alpha_\ell\) are compactly supported Lipschitz
functions. No bound for \(\nabla\alpha_\ell\) independent of
\(\eta\) is asserted or needed.
\end{lemma}

\begin{proof}
The tangent module is generated by gradients of Sobolev functions.
Thus finite sums of such gradients with bounded measurable scalar
coefficients are dense in \(L^2(TM)\). Approximating the coefficients
by finite-valued functions and refining the resulting partitions
shows that fields of the form
\[
 \sum_{\ell=1}^L\mathbf1_{E_\ell}\nabla q_\ell,
 \qquad q_\ell\in W^{1,2}(M),
\]
are dense as well: on each piece a constant linear combination of
gradients is the gradient of the corresponding linear combination
of potentials. This construction may be made simultaneously for
the finite list \(V_r\). Since \(f\) is bounded, it provides a
finite measurable partition and choices \(q_{r,\ell}\) for which
\[
 \sum_{r,\ell}\int_{E_\ell}
       f|\nabla q_{r,\ell}-V_r|^2<\eta/2.
\]

Put \(K_f=\operatorname{supp}f\). On this compact set the measure
with density
\[
 f\left(1+\sum_{r,\ell}|\nabla q_{r,\ell}-V_r|^2\right)
\]
is finite and regular. Choose disjoint compact subsets
\(K_\ell\subset E_\ell\cap K_f\) whose complement has arbitrarily
small measure for this density. There are nonnegative Lipschitz
functions \(\alpha_\ell\) summing to one, with
\(\alpha_\ell=1\) on \(K_\ell\). For example, take Lipschitz bumps
equal to one on the nonempty \(K_\ell\), supported in disjoint
neighborhoods, and assign their remaining complement to the first
function. On \(\bigcup_\ell K_\ell\), these functions agree with
the corresponding partition indicators. On the complement, the
change in the error integral is bounded by the sum appearing in
the displayed density. Taking the omitted measure sufficiently
small proves \eqref{fwd:selection-error}.
\end{proof}

\begin{proposition}[The curvature bound on compact test gradients]
\label{fwd:gradient-bound}
For every \(U,V\in\mathcal G\) and every nonnegative compactly
supported Lipschitz \(f\),
\begin{equation}
\label{fwd:core-bound}
 R(U,V,V,U)(f)\geq\kappa\int f|U\wedge V|^2.
\end{equation}
\end{proposition}

\begin{proof}
First retain an arbitrary finite sign-symmetric law as above. Apply
Lemma~\ref{fwd:gradient-selection} to the finite list
\[
 V_{AB}:=-\nabla_A(B-A).
\]
For each \(\ell\), apply \eqref{fwd:defective-curvature}, as
extended in Lemma~\ref{fwd:weight-approximation}, with weight
\(f\alpha_\ell\) and potentials \(q_{AB,\ell}\). Sum over
\(\ell\). The curvature functional and its model term are linear
in the weight, and \(\sum_\ell f\alpha_\ell=f\). Thus their sums
are exactly their evaluations at \(f\), including the terms
containing derivatives of the weight. The error is at most
\(\eta/2\), after decreasing the tolerance in
\eqref{fwd:selection-error} if needed. Letting \(\eta\downarrow0\)
gives
\begin{equation}
\label{fwd:averaged-bound}
 \mathbb ER(A,B,B,A)(f)
 \geq\kappa\int f\,\mathbb E|A\wedge B|^2.
\end{equation}
In particular there are no partition-derivative errors and no
uniform derivative bound on the partitions is needed.

Choose the uniform law on the four potentials yielding the fields
\(U,-U,V,-V\). In \eqref{fwd:averaged-bound}, pairs from the same
signed field have zero sectional and wedge expressions. The mixed
pairs have total probability \(1/2\). Multilinearity and the
curvature symmetries identify all of them with the corresponding
expression for \((U,V)\). Both sides of
\eqref{fwd:averaged-bound} are therefore one half of the respective
sides of \eqref{fwd:core-bound}. This proves the proposition.
\end{proof}

\subsection{The global test classes}

We supply the extension details, including the approximation
topology needed in the curvature formula. These details also explain
why the compact-gradient argument proves the exact global
quantifiers in Theorem~\ref{fwd:curvature}.

\begin{lemma}[Continuous coefficients on a fixed gradient family]
\label{fwd:coefficients}
Fix \(E_1,\ldots,E_N\in\mathcal G\). Then
\eqref{fwd:target} holds for nonnegative compact test weights and fields of the
form
\[
 X=\sum_{i=1}^N a_iE_i,
 \qquad Y=\sum_{i=1}^N b_iE_i,
 \qquad a_i,b_i\in\Test(M).
\]
\end{lemma}

\begin{proof}
For constant coefficient vectors \(a,b\in\mathbb R^N\), put
\(E(a)=\sum_i a_iE_i\). This is again in \(\mathcal G\).
Proposition~\ref{fwd:gradient-bound} says that
\[
 T_{a,b}(f):=
 R(E(a),E(b),E(b),E(a))(f)
 -\kappa\int f|E(a)\wedge E(b)|^2
\]
is nonnegative on nonnegative compact Lipschitz functions.
It is therefore a positive Radon measure. Indeed, for weights
supported in a fixed compact set, choose a nonnegative Lipschitz
cutoff dominating one there. Positivity bounds the functional by
its value on that cutoff times the uniform norm of the weight.
The functional extends to continuous compactly supported functions
and the Radon representation theorem applies.

Define functionals on compact Lipschitz weights by
\[
 \begin{gathered}
 D_{ijkl}=R(E_i,E_j,E_k,E_l)-\kappa Q_{ijkl}m,\\
 Q_{ijkl}=\langle E_i,E_l\rangle\langle E_j,E_k\rangle
          -\langle E_i,E_k\rangle\langle E_j,E_l\rangle.
 \end{gathered}
\]
For \(\alpha,\beta\in\mathbb N_0^N\) with
\(|\alpha|=|\beta|=2\), let
\[
 C_{\alpha\beta}
 :=\sum_{\substack{e_i+e_l=\alpha\\e_j+e_k=\beta}}D_{ijkl},
 \qquad
 T_{a,b}=\sum_{\alpha,\beta}a^\alpha b^\beta C_{\alpha\beta}.
\]
Here \(e_i\) are the coordinate vectors, and the first sum is
over ordered quadruples, retaining all monomial multiplicities.
Polynomial interpolation at finitely many constant vectors
expresses every grouped coefficient \(C_{\alpha\beta}\) as a
finite real linear combination of the positive Radon measures
\(T_{a,b}\). Thus each \(C_{\alpha\beta}\) is a signed Radon
measure; no measure representation of an individual \(D_{ijkl}\)
is needed.

On a fixed relatively compact region containing the support of
the weight, put
\[
 \lambda=\sum_{\alpha,\beta}|C_{\alpha\beta}|,\qquad
 c_{\alpha\beta}=\frac{dC_{\alpha\beta}}{d\lambda},\qquad
 P_x(a,b)=\sum_{\alpha,\beta}a^\alpha b^\beta
 c_{\alpha\beta}(x).
\]
For every rational pair \(a,b\), positivity of \(T_{a,b}\)
gives \(P_x(a,b)\geq0\) for \(\lambda\)-almost every \(x\).
Countability and polynomial continuity give a common conull set
on which this holds for every real pair.

For variable test coefficients, first apply tensoriality of
\(R\) over the test algebra. All scalar weights in the following
identity are compact test functions:
\[
 \begin{split}
 &R(X,Y,Y,X)(f)-\kappa\int f|X\wedge Y|^2\\
 &\qquad=\sum_{i,j,k,l}D_{ijkl}(fa_i a_l b_j b_k)\\
 &\qquad=\sum_{\alpha,\beta}C_{\alpha\beta}(fa^\alpha b^\beta)\\
 &\qquad=\int f(x)P_x(a(x),b(x))\dd\lambda(x)\geq0.
 \end{split}
\]
The last line uses only the continuous representatives of scalar
test functions on the signed measures. No tangent field is
evaluated on a singular measure support. This proves the claim,
including when some \(C_{\alpha\beta}\) are singular with
respect to \(m\).
\end{proof}

\begin{lemma}[Local approximation with strong second derivatives]
\label{fwd:field-approximation}
Let \(b\in\Test(M)\), and let \(E\subset M\) be compact. There are
compactly supported test functions \(b_j\) with bounded Laplacian
such that their gradients are uniformly bounded and, on a fixed
neighborhood of \(E\),
\[
 b_j\to b,\quad\nabla b_j\to\nabla b,\quad
 \Delta b_j\to\Delta b,\quad \Hess b_j\to\Hess b
\]
strongly in \(L^2\). The gradients converge almost everywhere
after passage to a subsequence.
Consequently, for any finite collection of test vector fields,
there are fields built with test coefficients and potentials of
this compact bounded-Laplacian class which converge locally to the
given fields, uniformly bounded in norm, with strong \(L^2\)
convergence of the fields, their covariant derivatives, and their
divergences.
\end{lemma}

\begin{proof}
Choose a fixed compact neighborhood \(E'\) of \(E\) and apply
Lemma~\ref{pre:approximation}\textup{(iii)} on \(E'\).
It gives the uniform gradient bound and all the asserted scalar
convergences on this fixed neighborhood. Strong \(L^2\)
convergence of the gradients gives almost-everywhere convergence
after passage to a subsequence.

For \(X=\sum_i a_i\nabla b_i\), keep the original coefficients
and approximate its finitely many potentials as just described.
The product rules give
\[
 \nabla X_j=\sum_i da_i\otimes db_{i,j}+a_i\Hess b_{i,j},
 \qquad
 \operatorname{div}X_j
 =\sum_i\langle\nabla a_i,\nabla b_{i,j}\rangle+a_i\Delta b_{i,j}.
\]
The coefficients and their gradients are bounded, so both
expressions converge strongly in \(L^2\) on the neighborhood.
The vector bounds and convergence follow from the corresponding
properties of the gradients. Use one common subsequence for the
finite collection.
\end{proof}

\begin{proof}[Proof of Theorem~\ref{fwd:curvature}]
It remains to extend Proposition~\ref{fwd:gradient-bound}; the RCD
and support conclusions were noted at the beginning of the section.
Fix a nonnegative compactly supported test weight \(f\) and
arbitrary \(X,Y\in\TestV(M)\). Lemma~\ref{fwd:field-approximation}
gives approximating fields \(X_j,Y_j\) near
\(\operatorname{supp}f\). Lemma~\ref{fwd:coefficients} proves
\eqref{fwd:target} for every pair \(X_j,Y_j\).

The field-continuity assertion of Lemma~\ref{pre:tensor-continuity}
passes the curvature evaluation to the limit on the fixed compact
support of \(f\). The wedge term converges by dominated convergence
on that support. This proves
\eqref{fwd:target} for arbitrary test fields and compact test weights.

Finally let \(f\in\Test(M)\) be nonnegative without a support
restriction. Choose \(0\leq\chi_R\leq1\) in the stronger test
cutoff class, equal to one on \(B_R(o)\), with compact support and
uniformly bounded gradient. Such an exhaustion exists for every
fixed \(K,n\) \cite[Appendix A, Lemma A.2]{GigliTamanini2021}.
The product \(f\chi_R\) is a compact test function, so the desired
inequality holds for this weight.

All the original tensor integrands are globally integrable. Test
fields are bounded and in \(L^2\), and their derivatives and
divergences are in \(L^2\). Terms with two derivatives are thus
\(L^1\); a typical term with \(df\) is bounded by a constant times
\(|W||\nabla Z|\), again in \(L^1\). The new terms containing
\(d\chi_R\) obey the same bound, with an additional indicator of
\(M\setminus B_R(o)\), and hence their integrals tend to zero.
The model term is integrable, since
\[
 f|X\wedge Y|^2
 \leq\|f\|_\infty\|X\|_\infty^2|Y|^2.
\]
Dominated convergence proves \eqref{fwd:target} for the original
global weight. This completes the proof.
\end{proof}

\begin{remark}[Order of limits]
\label{fwd:limit-order}
The construction uses successive limits, not a single uniform
expansion over its approximation parameters. First \(h\downarrow0\)
is taken for fixed fields, corrections, semiconcave weight, and
Sobolev dual potentials. Next the endpoint corrections approach
\(-\nabla_AA\) in \(L^2\). The defective inequality is then extended
to Lipschitz weights. Only after that extension are the partitions
and their associated dual potentials selected to remove the square.
The final field approximation is local on the support of a fixed
test weight, and the exhaustion is taken last. Thus neither the
divergence bounds of the endpoint corrections nor the gradients of
the localization partitions need to be uniform along their
approximating sequences.
\end{remark}
